% Printed 458; principal PDF 465; Google PDF 469.
\par\clearpage
\noindent\makebox[0pt][l]{\fontsize{13}{16}\selectfont 458}\hfill{\small\textbf{HÖHENMESSUNGEN.}}\hfill\null\par\medskip
\begin{center}
{\large Tafel III.}\par\medskip
\begingroup\setlength{\tabcolsep}{12pt}\renewcommand{\arraystretch}{1.05}
\begin{tabular}{r|r@{\hspace{12pt}\vrule width .9pt\hspace{12pt}}r|r@{\hspace{12pt}\vrule width .9pt\hspace{12pt}}r|r}
&$+$&&$+$&&$+$\\
1,9&1&2,8&4&3,4&17\\
2,3&1&2,9&5&3,5&22\\
2,4&2&3,0&7&3,6&27\\
2,5&2&3,1&9&3,7&34\\
2,6&3&3,2&11&3,8&43\\
2,7&3&3,3&14&3,9&54\\
2,8&4&3,4&17&&\\
\end{tabular}\endgroup
\par\bigskip\bigskip
{\large Gebrauch der Tafeln.}
\end{center}
\par\noindent\hspace*{\parindent}$t$, $t'$ Temperatur der Luft; $T$, $T'$ Temperatur des Quecksilbers (nach Réaumur); $b$, $b'$ Barometerstand (in \textls[120]{beliebiger} Maßeinheit).

\par\noindent\hspace*{\parindent}Man vermindert $\log b$ und $\log b'$ um $10T$ bzw. $10T'$ (als Einheiten der $5^{\mathrm{ten}}$ Dezimalstelle verstanden) und zieht die so korrigierten Logarithmen voneinander ab; die Differenz sei $=u$. Man addiert $\log u$ und $A$, nachdem man den letzteren Wert, falls man es für nötig hält, nach der zweiten Tafel korrigiert hat (sie gibt wie die dritte Tafel Einheiten in der $5^{\mathrm{ten}}$ Dezimalstelle an); die Summe sei $=v$. Diese Größe erhält noch eine kleine Korrektur aus Tafel III, deren Argument $v$ selbst ist. Das so korrigierte $v$ ist der Logarithmus des Höhenunterschieds in Metern. Will man den Höhenunterschied in Toisen haben, addiert man zum Logarithmus noch $9{,}71018$.
\par\medskip
\noindent\hspace*{.5em}\begin{tabular}{l@{\quad}l@{\quad}l@{\quad}l@{\quad}l}
Beispiel.&$t=15^\circ{,}3,$&$T=14{,}9,$&$b=735{,}581\,\mathrm{mm},$&Polhöhe $=45^\circ$.\\
&$t'=3{,}2,$&$T'=7{,}8,$&$b'=537{,}203$&
\end{tabular}
\par\noindent\hspace*{1.5cm}\begingroup\setlength{\tabcolsep}{2pt}\renewcommand{\arraystretch}{1.2}
\begin{tabular}{r@{\ }c@{\ }r@{\quad}c@{\ }r@{\ }l}
$\log b$&$=$&$2{,}86663$&Korr.&$-149$&\\
$\log b'$&$=$&$2{,}73014$&»&$-78$&\\\cline{2-3}\cline{5-5}
&&$0{,}13649$&»&$-71$&\qquad .\\
$u$&$=$&$0{,}13578$&&&\\\cline{2-3}
$\log u$&$=$&$9{,}13284$&&&\\
$A$&$=$&$4{,}28408$&&&\\
Korr.&$=$&$0$&&&\\\cline{2-3}
$v$&$=$&$3{,}41692$&Korr.&$+18$&$.\ .\ .\ .\ 3{,}41710=\log2612{,}8\,\mathrm{m}.$
\end{tabular}\endgroup
\begin{center}\rule{3.5cm}{.4pt}\end{center}

% Unnumbered title [459]; principal PDF 466; Google PDF 470.
\clearpage\null\vspace*{6.7cm}
\begin{center}{\fontsize{24}{30}\selectfont\textls[160]{HELIOTROP.}}\end{center}
\par\vspace*{11cm}\noindent\makebox[.86\linewidth][r]{58\textsuperscript{*}}\par
% Blank reverse [460]; principal PDF 467; Google PDF 471.
\clearpage\null
% Unnumbered opening [461]; principal PDF 468; Google PDF 472.
\clearpage\vspace*{2.7cm}
\begin{center}
\begin{tabular}{c}\hline
\small\textbf{Göttingische gelehrte Anzeigen.\quad 126. Stück.\quad 9. August 1821.\quad S. 1249--1254.}\\\hline
\end{tabular}\par\bigskip
{\large [Über den Heliotrop.]}
\end{center}
\par\noindent\hspace*{\parindent}Wer mit der höheren Geodäsie vertraut ist, kennt die Schwierigkeiten, für die Bildung \textls[120]{großer} Dreiecke wirklich geeignete Zielpunkte zu finden. Hohe Kirchtürme stehen in manchen Gegenden nicht an geeigneten Stellen, und auch die vorhandenen Türme bieten oft nicht die gewünschten Möglichkeiten, die Instrumente aufzustellen und die gemessenen Winkel zu zentrieren. Auch sind sie häufig nicht so regelmäßig gebaut, wie es für die größtmögliche Genauigkeit wünschenswert ist. Eigens gebaute Signaltürme haben, abgesehen von dem Aufwand an Geld und Zeit für ihre Errichtung, mit den Kirchtürmen gemeinsam, dass sie sich aus beträchtlicher Entfernung schwer erkennen und anvisieren lassen, wenn sie sich vor einem nahen dunklen Hintergrund abzeichnen. Gibt man ihnen eine helle Farbe, zeigen sie je nach Sonnenbeleuchtung eine wechselnde und äußerst nachteilige Phase. Selbst die vollkommensten Signaltürme, die geschwärzt sind und sich gegen den Himmel abzeichnen, sind in \textls[120]{sehr} großen Entfernungen nicht ganz frei von einer störenden Phase, wenn man zugleich eine von der Sonne beleuchtete und eine im Schatten liegende Seitenfläche sieht. Messungen bei Nacht mithilfe von \textsc{Argand}-Lampen sind diesen Fehlern zwar nicht unterworfen, bringen dafür aber, besonders auf schwer zugänglichen Bergen, andere Unannehmlichkeiten mit sich, die zu offensichtlich sind, als dass sie hier besprochen werden müssten.

\par\noindent\hspace*{\parindent}Diese Überlegungen haben Herrn Hofrat \textsc{Gauss} veranlasst, für die

% Printed 462; principal PDF 469; Google PDF 473.
\par\clearpage
\noindent\makebox[0pt][l]{\fontsize{13}{16}\selectfont 462}\hfill{\small\textbf{HELIOTROP.}}\hfill\null\par\medskip
\noindent auf allerhöchsten Befehl im Königreich Hannover durchzuführende Gradmessung ein neues Hilfsmittel zu entwerfen, das sich, und sei es nur neben dem genannten, bei der großräumigen Triangulation vorteilhaft einsetzen ließe. Der Erfolg hat seine Erwartungen noch weit übertroffen.

\par\noindent\hspace*{\parindent}Eine auf photometrische Überlegungen gestützte Untersuchung hatte ihn schon früher davon überzeugt, dass das von einem sehr kleinen Planspiegel reflektierte Sonnenlicht selbst in den größten bei Triangulationen vorkommenden Entfernungen noch stark genug sein müsse, um einen ausgezeichneten Zielpunkt abzugeben. Um diese Idee zu nutzen, kam es darauf an, ein Instrument zu entwerfen, mit dessen Hilfe das Sonnenlicht überall genau in jede nötige Richtung gelenkt werden kann. Zugleich musste sich ein solches Instrument überall leicht aufstellen und bedienen lassen, und der Mittelpunkt des reflektierenden Spiegels musste bei den Bewegungen, die nötig sind, um der wandernden Sonne gleichsam zu folgen, immer vollkommen in Ruhe bleiben. Für ein solches Instrument, das diese \textls[120]{Lenkung des Sonnenlichts} in jede beliebige Richtung auf die beschriebene Weise vollkommen auszuführen vermag, scheint der Name \textls[120]{Heliotrop} am besten geeignet, mindestens ebenso gut wie für zwei bekannte Produkte des Pflanzen- und Mineralreichs.

\par\noindent\hspace*{\parindent}Weitere Überlegungen zu diesem Gegenstand haben Herrn Hofrat \textsc{Gauss} zu zwei ganz unterschiedlichen Bauweisen eines Heliotrops geführt. Unser geschickter Herr Inspektor \textsc{Rumpf} hat bereits ein solches Instrument nach der einen Bauweise vorzüglich ausgeführt, und an einem zweiten nach der anderen Bauweise arbeitet derselbe Handwerker jetzt. Eine vollständige Beschreibung, die für unsere Blätter ungeeignet wäre, wird an anderer Stelle gegeben. --- Um noch vor der Fertigstellung des genannten Heliotrops wirkliche Versuche über die Stärke des reflektierten Sonnenlichts machen zu können, kam Herr Hofrat \textsc{Gauss} noch auf eine dritte Idee. Mit ihr lässt sich jeder Spiegelsextant zu einem recht vollkommenen Heliotrop umbauen, ja nötigenfalls sogar ohne jeden Zusatz als Heliotrop verwenden, wenn er nur auf ein gutes Stativ gesetzt werden kann. Herr Hofrat \textsc{Gauss} ließ nämlich Herrn Inspektor \textsc{Rumpf} am großen Spiegel eines Sextanten einen dritten Planspiegel so befestigen, dass dessen

% Printed 463; principal PDF 470; Google PDF 474.
\par\clearpage
\noindent\null\hfill{\small\textbf{HELIOTROP.}}\hfill\makebox[0pt][r]{\fontsize{13}{16}\selectfont 463}\par\medskip
\noindent Ebene senkrecht auf der Ebene des Sextanten steht und mit der Ebene des großen Spiegels einen Winkel bildet, der den Winkel zwischen der Visierlinie (die nötigenfalls erst durch Einziehen feiner Kreuzfäden festzulegen ist) und der Ebene des kleinen Spiegels zu einem rechten Winkel ergänzt. Sobald man einen solchen Sextanten so ausrichtet, als wollte man den Winkelabstand eines Objekts vom Mittelpunkt (ja, sofern alles gut gearbeitet und justiert ist, von irgendeinem Punkt) der Sonnenscheibe messen, reflektiert der dritte Spiegel das Sonnenlicht in Richtung des Objekts, gleichgültig, welches von beiden man direkt sieht. Steht ein solcher Sextant auf einem guten Stativ, lässt sich mit etwas Übung das reflektierte Sonnenlicht ohne Unterbrechung zum gewünschten Punkt senden. Fehlt der dritte Spiegel, kann der große Spiegel selbst an seine Stelle treten, wenn man bei vollkommen unveränderter Ebene die Alhidade rasch um den zuvor genannten Winkel vorschiebt (oder nominell auf dem Gradbogen um den doppelten Winkel). Mit dieser letzten Methode lässt sich offensichtlich nur eine unterbrochene Reflexion erzielen. Doch kann man mit Übung, der nötigen Sorgfalt und ein paar kleinen Kunstgriffen, deren Beschreibung hier zu weit führen würde, das Sonnenlicht jeweils zwei Minuten und länger zum vorgegebenen Punkt reflektieren. Die beiden letzteren Methoden haben übrigens offensichtlich den kleinen Nachteil, dass der Mittelpunkt des reflektierenden Spiegels bei feststehendem Stativ nicht vollkommen in Ruhe bleibt. In den meisten Fällen wird dies jedoch so gut wie unerheblich sein. Hält man es für nötig, kann man auch leicht fortwährend etwas nachjustieren oder dies bei der Berechnung berücksichtigen. Die Anwendbarkeit ist natürlich auch auf den durch die Größe des Gradbogens am Sextanten vorgegebenen Winkelabstand der Sonne vom Objekt beschränkt.

\par\noindent\hspace*{\parindent}Bevor wir die Ergebnisse der Versuche mit dem Heliotrop anführen, sei bemerkt, dass alle zur Reflexion verwendeten Spiegel 2 Zoll breit und $1\frac14$ Zoll hoch sind. Die Erfahrung hat bestätigt, was Herr Hofrat \textsc{Gauss} schon aus photometrischen Gründen vorausberechnet hatte, nämlich dass bei einigermaßen günstigen Umständen größere Abmessungen ganz unnötig wären, wenigstens für den geodätischen Gebrauch. Bei der zweiten oben erwähnten Bauweise kann man übrigens
% Printed 464; principal PDF 471; Google PDF 475.
\par\clearpage
\noindent\makebox[0pt][l]{\fontsize{13}{16}\selectfont 464}\hfill{\small\textbf{HELIOTROP.}}\hfill\null\par\medskip
\noindent auch nach Belieben einen größeren Spiegel anbringen lassen, ohne sonst die Abmessungen des Instruments zu vergrößern.
\begin{center}
\begin{tikzpicture}[x=.007cm,y=-.007cm,line cap=round,line join=round]
\path[use as bounding box] (0,0) rectangle (2030,910);
\draw[line width=.65pt] (655,75) circle[radius=.238cm];
\fill (655,75) circle[radius=.0301cm];
\draw[line width=.45pt] (702,125)--(1165,750);
\draw[line width=.45pt] (78,750)--(1165,750)--(1976,740);
\draw[line width=.65pt] (644,735)--(977,735)--(979,740)--(979,760)--(975,765)--(646,766)--(641,760)--(641,741)--cycle;
\draw[line width=.65pt] (608,743)--(641,741)--(642,759)--(610,759)--(606,754)--cycle;
\draw[line width=.5pt] (877,618)--(1413,858);
\draw[line width=.5pt] (1234,607)--(1102,868);
\draw[line width=.5pt,dash pattern=on 0pt off 2.7pt] (884,634)--(1403,867);
\draw[line width=.5pt,dash pattern=on 0pt off 2.7pt] (1239,618)--(1116,867);
\end{tikzpicture}
\end{center}
\par\noindent\hspace*{\parindent}Über die Entfernung von der hiesigen Sternwarte zum Hohehagen, einem Hauptdreieckspunkt der Gradmessung (beinahe 2 geografische Meilen), war das Licht sowohl des eigentlichen Heliotrops als auch des zum Heliotrop umgebauten Sextanten bei hellem Sonnenschein mit bloßem Auge überaus schön zu sehen. Im Fernrohr des Theodoliten war es eigentlich zu stark, während allein das von einer hellen Wolke reflektierte Licht den schönsten denkbaren Zielpunkt ergab. Offensichtlich lässt sich das reflektierte Sonnenlicht selbst, wo man es wünscht, leicht nach Belieben abschwächen, indem man einen Teil des Spiegels abdeckt.

\par\noindent\hspace*{\parindent}Über die Entfernung vom Hils (einem anderen Hauptdreieckspunkt) zum Meridianzeichen der Sternwarte, fast 5 geografische Meilen, war das Licht beider Instrumente ebenfalls noch mit bloßem Auge ausgezeichnet zu sehen, wie ein schönes Sternchen, und bot im Fernrohr des Theodoliten einen hervorragenden Zielpunkt. Bei nebliger Luft ließ sich manchmal weder vom Bergrücken des Hils noch vom dort gebauten Signalturm im Fernrohr des Theodoliten eine Spur erkennen. Das Heliotroplicht schien dann wie ein prachtvoller Stern am blauen Himmel zu schweben.

\par\noindent\hspace*{\parindent}Die wichtigsten Versuche wurden erst in den letzten Tagen angestellt. Herr Professor \textsc{Encke}, Leiter der Seeberger Sternwarte, war auf Einladung von Herrn Hofrat \textsc{Gauss} hierhergekommen, um den Gebrauch des Spiegelsextanten ohne dritten Spiegel als Heliotrop und die dabei anzuwendenden

% Printed 465; principal PDF 472; Google PDF 476.
\par\clearpage
\noindent\null\hfill{\small\textbf{HELIOTROP.}}\hfill\makebox[0pt][r]{\fontsize{13}{16}\selectfont 465}\par\medskip
\noindent kleinen Kunstgriffe kennenzulernen. Er begab sich dann auf den Inselsberg, während Herr Hofrat \textsc{Gauss} mit den Messungen auf dem Hohehagen begann. Ersterer sandte das Sonnenlicht mit dem als Heliotrop verwendeten Sextanten in Abständen zum Hohehagen (Entfernung 85000 Meter oder $11\frac12$ geografische Meilen), von wo es mit dem eigentlichen Heliotrop zum Inselsberg gelenkt wurde. Die Versuche und Beobachtungen wurden vom 19. bis 29. Juli unter abwechselnd ungünstigen und günstigen Bedingungen fortgesetzt und hatten den bestmöglichen Erfolg. Das Heliotroplicht lieferte beiden Beobachtern die schönsten denkbaren Zielpunkte. Häufig erschien es wie ein schönes Sternchen, während man im selben Fernrohr den Umriss des Berges kaum oder gar nicht erkennen konnte. Einer der Beobachter befand sich zuweilen in Nebel und Regen, während das Heliotroplicht von der anderen Seite kräftig hindurchdrang. Einige Male glaubten sogar mehrere Anwesende mit besonders scharfen Augen auf dem Hohehagen, den Lichtpunkt auf dem Inselsberg mit bloßem Auge erkennen zu können. Wir können noch hinzufügen, dass die auf das Heliotroplicht bezogenen Winkelmessungen selbst auf beiden Seiten so gut übereinstimmten, wie man es in einer so großen Entfernung von keinem anderen Signal außer unter ganz besonders günstigen Umständen hätte erwarten dürfen.

\par\noindent\hspace*{\parindent}Diese Erfahrungen zeigen bereits zweifelsfrei, dass bei der Verwendung von Heliotroplicht der Größe der zu bildenden Dreiecke nur noch die Grenzen gesetzt sind, die sich aus der Erdkrümmung ergeben.

\par\noindent\hspace*{\parindent}Der Bedarf der höheren Geodäsie hat den Anstoß zu diesem Instrument gegeben. Deshalb beschränken wir uns hier darauf, die genannten Erfahrungen zu schildern, ohne die naheliegende Aussicht auf einen künftigen, vielleicht noch wichtigeren Gebrauch eines so kräftig den Raum durchdringenden Mittels für die telegrafische Signalgebung in Krieg und Frieden jetzt weiterzuverfolgen.
\par\vfill\noindent\makebox[0pt][l]{IX.}\makebox[.86\linewidth][r]{59}\par

% Printed 466; principal PDF 473; Google PDF 477.
\clearpage
\noindent\makebox[0pt][l]{\fontsize{13}{16}\selectfont 466}\hfill{\small\textbf{HELIOTROP.}}\hfill\null\par\bigskip\bigskip
\begin{center}
\begin{tabular}{c}\hline
\small\textbf{Astronomisches Jahrbuch für das Jahr 1825.\quad Berlin 1822.\quad S. 103 und 104.}\\\hline
\end{tabular}\par\bigskip
{\large Erfindung eines Heliotrops,\\
von Herrn Hofrat Ritter \textsc{Gauss} in Göttingen unter dem 26. Dezember 1821\\
eingesandt.}
\end{center}
\par\noindent\hspace*{\parindent}Sie werden mein langes Schweigen entschuldigen, wenn ich Ihnen sage, dass ich den größten Teil des Jahres nicht hier war und sogar noch im Spätherbst nach Altona gereist bin, um den \textsc{Ramsden}schen Zenitsektor in Empfang zu nehmen. Von dort bin ich erst vor Kurzem hierher zurückgekehrt.

\par\noindent\hspace*{\parindent}Bei meiner Triangulation, bei der ich bisher an fünf Dreieckspunkten die Winkel gemessen habe, habe ich versucht, die Dreiecke so groß wie möglich zu machen. Über das neue Hilfsmittel, das ich hierfür eingesetzt habe, den \textls[120]{Heliotrop}, und über die ersten großräumigen Versuche damit werden Sie die Nachricht in Nr. 126 der hiesigen gelehrten Anzeigen gelesen haben. Seitdem habe ich ihn ständig verwendet, sowohl als Zielpunkt bei Winkelmessungen als auch mit nicht weniger Erfolg zur telegrafischen Signalgebung. Die gewaltige Wirkung des von einem Spiegel von 2 Zoll Breite und $1\frac14$ Zoll Höhe reflektierten Sonnenlichts, das in Entfernungen von 5, 6, $7\frac12$, ja einmal sogar von $9\frac13$ geografischen Meilen mit \textls[120]{bloßem Auge} gesehen wurde, versetzt gewöhnlich diejenigen in Erstaunen, die es zum ersten Mal erleben und nicht durch theoretische Berechnung darauf vorbereitet sind. Bei einigermaßen günstiger Beschaffenheit der Luft sind der Größe der Dreiecksseiten jetzt nur noch durch die Erdkrümmung Grenzen gesetzt, zumal wenn man die Spiegel etwas größer macht, wie ich es bei zwei neu angefertigten Heliotropen ganz verschiedener Bauweise getan habe.
\begin{center}\rule{3.5cm}{.4pt}\end{center}

% Printed 467; principal PDF 474; Google PDF 478.
\par\clearpage
\noindent\null\hfill{\small\textbf{BRIEFWECHSEL.\quad HELIOTROP.}}\hfill\makebox[0pt][r]{\fontsize{13}{16}\selectfont 467}\par
\vspace*{4.2cm}
\begin{center}{\fontsize{19}{24}\selectfont\textls[80]{BRIEFWECHSEL.}}\end{center}
\vspace*{1.5cm}
\begin{center}{\large\textsc{Gauss} an \textsc{Olbers}.\quad Göttingen, 1. Juli 1821.}\end{center}
\par\noindent\hspace*{\parindent}.\ .\ .\ .\ .\ .\enspace Mir kam es inzwischen zunächst nur darauf an, einige Erfahrungen über die Strahlkraft der Spiegel selbst zu sammeln. Zuerst habe ich Verschiedenes versucht. Ich befestigte einen Spiegel am Deckel des Fernrohrs eines Theodoliten und versuchte, ihn anhand zuvor mühsam berechneter Azimute und Höhen so auszurichten, dass er das Licht in eine bestimmte Richtung warf. Das misslang jedoch völlig. Der Deckel war etwas schwergängig und ließ sich nicht zuverlässig immer wieder in dieselbe Stellung bringen. Es blieben Abweichungen von $20'$, die dieses Verfahren ganz unbrauchbar machten, sofern der Spiegel nicht fester am Fernrohr befestigt werden konnte, sodass dieses offen blieb. Der Ärger über die vergebliche Mühe brachte mich inzwischen auf eine andere Idee, die sich vollkommen bewährt hat. Schon der bloße Spiegelsextant auf einem guten Stativ leistet das Verlangte, wenn auch nicht so vollkommen wie ein eigentlicher Sonnenspiegel. Beträgt der Winkel, den die Visierlinie (die, sofern sie nicht schon vorhanden ist, erst durch einen Faden oder ein Fadenkreuz festgelegt werden muss) mit dem kleinen Spiegel bildet, $=90^\circ-a$, und sind das Sonnenbild und das Objekt, zu dem das Licht geworfen werden soll, wie üblich zur Berührung gebracht, als wollte man den Winkel messen --- gleichgültig, ob man ersteres oder letzteres direkt sieht ---, braucht man nur bei unveränderter Ebene die Alhidade um $a$ (oder nominell $2a$) vorzuschieben, um das Gewünschte zu erreichen. Mit etwas Übung lässt sich die Stellung leicht so wählen, dass jene Koinzidenz
\par\medskip\noindent\makebox[.86\linewidth][r]{59\textsuperscript{*}}\par
% Continuation on printed 468.
